
\begin{frame}{TrustZone / Secure World}
  \begin{itemize}
    \item \href{https://developer.arm.com/documentation/100935/0100/The-TrustZone-hardware-architecture-?lang=en}{TrustZone} is an ARM-only hardware isolation mechanism  
    \item It is implemented on the \href{https://developer.arm.com/Architectures/A-Profile\%20Architecture}{Application(A)-profiles} of ARMv7-9, this is what we discuss
    \item Also implemented on the ARMv7-M (Microcontroller) profile, but differently
    \item Splits the system into 2 \textbf{worlds}: Secure and non-Secure
    \item They are also called \textbf{Trusted} / \textbf{Rich} Execution Environment (TEE/REE)
    \item This is \textbf{orthogonal} to Exception Levels:
    \begin{itemize}
      %% Technically, Chapter D1.1 the Arm Architecture Reference Manual for A profile Architecture
      %% ()
      %% lists both EL2 and EL3 as optional unless on ARMv9 with the Realms Management Extenstion.
      %% In practice, buying an A profile core without EL3 is probably not worth it.
      \item EL3 (Secure Monitor) only exists in \textbf{Secure} world
      \item EL2, if it exists, can exist in \textbf{Secure}, \textbf{Non-Secure}, or both
      %% This is not accurate for ARMv9 with RME (Realms Mangement Extensions)
      \item EL1 and EL0 must exist in \textbf{both} security states
    \end{itemize}
  \end{itemize}
\end{frame}

\begin{frame}{TrustZone / Secure World}
  \center\includegraphics[width=0.7\textwidth]{slides/security-hw-sec-barriers/Exception_Levels.pdf}
\end{frame}

\begin{frame}{World transitions: Secure to Non-Secure}
  \begin{itemize}
    \item The CPU starts executing at the highest EL: EL3, so in Secure World
    \item It then sets up the rest of Secure World, then Normal World, in stages
    \item Similar to the EL, the CPU keeps track of the security state in a register
    \begin{itemize}
      \item SCR is the \textit{Secure Configuration Register (SCR)}, bit 0 is the \textit{NS} bit
      \item SCR can only be accessed in Secure EL1-3
      \item Going from Secure World back to Non-Secure is simply setting \textit{NS} to 1
      \item The Security model expects only the Secure Monitor to do this transition
    \end{itemize}
    \item The ARM spec leaves non-EL3 access to SCR as \textit{undefined}
  \end{itemize}
\end{frame}

\begin{frame}{World transitions: Back to Secure}
  \begin{itemize}
    \item The systems transitions from Non-Secure to Secure via exceptions, either:
    \begin{itemize}
      \item asynchronous, via some \textit{FIQ} interrupts, depending on \textit{SCR.FIQ}
      \item synchronous, generated by an \textit{smc} (\textit{secure monitor call}) instruction
    \end{itemize}
    \item \textit{smc} cannot be called from EL0, so userland would first need to issue a \textit{svc}
    \item The authority is the \textit{Secure Configuration Register (SCR)}'s NS bit
    \item Accessing SCR from non-secure or from S-EL0 will fault.
    \item Only the Secure Monitor (running at EL3) \textbf{should} change it.
  \end{itemize}
\end{frame}

\begin{frame}{\href{https://www.trustedfirmware.org/}{Arm Trusted Firmware}}
  \begin{itemize}
    \item Trusted Firmware is an open-source reference Secure World implementation
    \item It is the default Secure World for most ARM platforms
    \item Comes in two flavours:
    \begin{itemize}
      \item Trusted Firmware A (for the Application profile), or
            \href{https://www.trustedfirmware.org/projects/tf-a}{TF-A}
      \item Trusted Firmware M (for the Microcontroler profile)
            \href{https://www.trustedfirmware.org/projects/tf-m}{TF-M}
    \end{itemize}
    \item There is now a Rust implementation of TF-A called Rusted Firmware A
          (\href{https://www.trustedfirmware.org/projects/rf-a}{RF-A})
    \item Most importantly, TF-A implements the Secure Monitor
    \item Uses a specific nomenclature to designate boot components
  \end{itemize}
\end{frame}

\begin{frame}{Trusted Firmware A (TF-A)}
  \begin{itemize}
    \item TF-A is split into 5 steps:
    \begin{itemize}
      \item BootLoader 1, or BL1: the SoC vendor's ROM code
      \item BootLoader 2, or BL2: "Trusted Boot firmware"
      \item BootLoader 3-1, or \href{https://elixir.bootlin.com/arm-trusted-firmware/v2.14.0/source/bl31}{BL31}:
            EL3 Runtime
      \item BootLoader 3-2, or \href{https://elixir.bootlin.com/arm-trusted-firmware/v2.14.0/source/bl32}{BL32} (optional):
            Secure-EL1 payload (Trusted OS, Trusted kernel, TEE)
      \item BootLoader 3-3, or BL33: Non-trusted firmware, the "normal" bootloader
    \end{itemize}
    \item Only BL31 and BL32 are resident after boot is complete:
    \begin{itemize}
      \item BL31 is responsible for routing Secure exceptions (and smcs) to the Secure World
      \item BL32 handles some of these exceptions and smcs, potentially on behalf of TAs
    \end{itemize}
  \end{itemize}
\end{frame}

\begin{frame}{ARMv8 boot sequence: TF-A names}
\includegraphics[width=0.7\textwidth]{common/armv8-boot-sequence-generic.pdf}
\end{frame}

\begin{frame}{Trusted Firmware A (TF-A)}
  \begin{itemize}
    \item In a "standard" boot, where U-Boot is used as boot loader:
    \begin{itemize}
      \item BL1: the vendor bootROM, non-resident
      \item BL2: is implemented by the \href{https://docs.u-boot.org/en/latest/usage/spl_boot.html}{Second Program Loader} (SPL)
      \item BL31: the runtime component of TF-A resident in EL3
      \item BL32: OP-TEE
      \item BL33: U-Boot proper
    \end{itemize}
  \end{itemize}
\end{frame}

\begin{frame}{ARMv8 boot sequence}
\includegraphics[width=0.7\textwidth]{common/armv8-boot-sequence-common-names.pdf}
\end{frame}

\begin{frame}{Rationale}
  \begin{itemize}
    \item The aim of TrustZone is to create a Trusted Execution Environment (TEE)
    \item The isolation is backed by the hardware
    \item Some hardware resources are only accessible to the TEE
    \item Sounds familiar, why is it not simply the kernel?
    \begin{itemize}
      \item attack surface: the TEE should be as small as possible
      \item usage: can you ensure attackers do not get \textit{root}
      \item virtualization: offer services securely to kernels you don't trust
    \end{itemize}
    \item The key concept is defense in depth
  \end{itemize}
\end{frame}

\begin{frame}{Trusted Execution Environments}
  \begin{itemize}
    \item The TEE is the part of Secure World that does things for the user
    \item This is an A profile concept, the M profile equivalent is Secure Processing Environment (SPE)
    \item There are different implementations of TEEs:
    \begin{itemize}
      \item Trustsonic's \href{https://www.trustonic.com/secure-os/technology/}{Kinibi}
      \item Samsung's \href{https://www.samsungknox.com/en}{Knox}
      \item Qualcomm's \href{https://docs.qualcomm.com/doc/80-88500-4/topic/78_Qualcomm_TEE_5_0.html}{QTEE}
      \item Huawei's iTrustee
    \end{itemize}
    \item Some are even Open Source:
    \begin{itemize}
      \item Android's \href{https://source.android.com/docs/security/features/trusty}{Trusty},
            based on Little Kernel (LK)
      \item Nvidia's Trusted Little Kernel (TLK), also based on LK: \code{git://nv-tegra.nvidia.com/3rdparty/ote_partner/tlk.git}
      \item OP-TEE, the reference Open Source TEE
    \end{itemize}
  \end{itemize}
\end{frame}

\begin{frame}{Trusted Execution Environments}
  \begin{itemize}
    \item The TEE is an optional component
    \item It does not normally offer service to Normal World directly
    \item It does provide HW access to Trusted Applications
    \item The TEE is the Secure World equivalent of the kernel, running in S-EL1
    \item For further isolation, ARMv8 introduces Secure Partitions
    \begin{itemize}
      \item Essentially run different instances of TEEs in isolated HW
      \item A Secure Partition Manager runs at S-EL2 (SPMC) and (S-)EL3 (SPMD)
    \end{itemize}
    \item ARMv9 pushed this even further with the Realms extension
  \end{itemize}
\end{frame}

\begin{frame}{Trusted Applications}
  \begin{itemize}
    \item Trusted Applications are the secure world "userland"
    \item OP-TEE has 2 main types of Trusted Apps:
    \begin{itemize}
      \item Pseudo-TAs, statically compiled into the OP-TEE core binary
      \item Usermode TAs are dynamically loaded by OP-TEE and run at S-EL0.\\
    \end{itemize}
    \item Usermode TAs can be further subdivided into:
    \begin{itemize}
      \item Early TAs: still stored in OP-TEE core, in a \code{data} section
      \item REE FS TAs: loaded by the REE via \code{tee-supplicant}.
    \end{itemize}
    \item Because REE FS TAs are loaded from the REE, the TEE needs to authenticate them
    \begin{itemize}
      \item OP-TEE expects REE FS TAs to be \href{}{signed}
      \item the OP-TEE core embarks the public part of
            \href{https://elixir.bootlin.com/op-tee/4.9.0/source/keys/default_ta.pem}{\code{default_ta.pem}}
      \item do \textbf{NOT} ship OP-TEE with that key
    \end{itemize}
  \end{itemize}
\end{frame}

\begin{frame}{Usecases}
  \begin{itemize}
    \item "Software HSM": this is why OP-TEE has a \href{https://github.com/OP-TEE/optee_os/tree/master/ta/pkcs11}{PKCS\#11 TA}
    \item \href{https://github.com/OP-TEE/optee_ftpm}{fTPM}: OP-TEE 
    \item Payment processing
    \item DRM
    \item Biometrics (fingerprint unlock)
    \item Hardware-isolated detection of non-secure malware
  \end{itemize}
\end{frame}

\begin{frame}{Issues}
  \begin{itemize}
    \item TrustZone is, ultimately, another security layer
    \item The security guarantees it offers are only as strong as the implementation
    \item Although it is hardware-backed, it is only useful if software is involved
    \item Like all software, Secure World software can have bugs
    \begin{itemize}
      \item \href{https://blog.thalium.re/posts/pivoting_to_the_secure_world/}{Arbitrary code execution in a Samsung TEE}
      \item \href{https://www.keysight.com/blogs/en/authors/federico-menarini}{Other TEEGris vulnerabilities}
      \item \href{https://blog.quarkslab.com/attacking-the-arms-trustzone.html}{S-EL0 stack buffer overflow in Qualcomm's TEE}
      \item \href{https://blog.quarkslab.com/a-deep-dive-into-samsungs-trustzone-part-3.html}{EL3 execution via Samsung Kinibi}
    \end{itemize}
    \item Secure world configuration can be very HW-dependent
  \end{itemize}
\end{frame}
